% Part: model-theory
% Chapter: interpolation
% Section: interpolation-proof

\documentclass[../../../include/open-logic-section]{subfiles}

\begin{document}

\olfileid{mod}{int}{prf}

\olsection{క్రేగ్ అంతర్వేశ సిద్ధాంతం}

\begin{thm}[క్రేగ్ అంతర్వేశ సిద్ధాంతం]
\ollabel{thm:interpol}
$\Entails !A \lif !B$ అయితే, $\Entails !A \lif !C$,
$\Entails !C \lif !B$ అయ్యే \tetoken{ఒక వాక్యం}{!!a{sentence}}~$!C$
ఉంటుంది; ఇంకా $!C$లోని ప్రతి
\tetoken{స్థిరసంకేతం}{!!{constant}},
\tetoken{ప్రమేయ సంకేతం}{!!{function}},
\tetoken{విధేయ సంకేతం}{!!{predicate}} ($\eq$ తప్ప) $!A$, $!B$
రెండింటిలోనూ వస్తుంది. ఆ \tetoken{వాక్యం}{!!{sentence}}~$!C$ను
$!A$, $!B$ల \emph{అంతర్వేశకం (ఇంటర్‌పోలంట్)} అంటారు.
\end{thm}

\begin{proof}
$!A$ యొక్క భాష $\Lang{L}_1$, $!B$ యొక్క భాష $\Lang{L}_2$
అనుకుందాం. $\Lang{L}_0 = \Lang{L}_1 \cap \Lang{L}_2$గా ఉంచండి.
ప్రతి $i \in \{0, 1, 2 \}$కు, అనంత సంఖ్యలో కొత్త
\tetoken{స్థిరసంకేతాలు}{!!{constant}s}~$\Obj c_0, \Obj c_1,
\Obj c_2, \dots$ను $\Lang{L}_i$కు చేర్చి $\Lang{L}'_i$ను పొందండి.

$!A$ అసంతృప్తిపరచదగినదైతే $\lexists[x][\eq/[x][x]]$ ఒక
అంతర్వేశకం. $\lnot !B$ అసంతృప్తిపరచదగినదైతే (అందువల్ల $!B$
చెల్లుబాటు అయితే) $\lexists[x][\eq[x][x]]$ ఒక అంతర్వేశకం. కాబట్టి
$!A$, $\lnot !B$ రెండూ సంతృప్తిపరచదగినవని కూడా అనుకోవచ్చు.

అంతర్వేశ సిద్ధాంతపు విపర్యయాన్ని నిరూపించడానికి $!A$, $!B$లకు
అంతర్వేశకం లేదని అనుకుందాం. మరో మాటలో చెప్పాలంటే,
$\Lang{L}_0$లో $\{!A\}$, $\{\lnot !B\}$లు వేరుపరచలేనివని
అనుకుందాం.

$(\{ !A \}, \{\lnot!B\})$ అనే జతను గరిష్ఠంగా వేరుపరచలేని జత
$(\Gamma^*, \Delta^*)$గా విస్తరించడమే మన లక్ష్యం. $\Lang{L}'_1$లోని
\tetoken{వాక్యాలను}{!!{sentence}s} $!A_0$, $!A_1$, $!A_2$, \dots{}గా
జాబితా చేద్దాం; $\Lang{L}'_2$లోని
\tetoken{వాక్యాలను}{!!{sentence}s} $!B_0$, $!B_1$, $!B_2$, \dots{}గా
జాబితా చేద్దాం.
\sourcecorrection{OLTEMODINT-003}{నిరూపణ చివర్లో
$\Gamma^*$ను $\Lang{L}'_1$లోను, $\Delta^*$ను $\Lang{L}'_2$లోను
గరిష్ఠ అవైరుధ్య సమితులుగా ఉపయోగిస్తుంది. కానీ ఆంగ్ల మూలం ఇక్కడ
విస్తరించని $\Lang{L}_1$, $\Lang{L}_2$ వాక్యాలను మాత్రమే లెక్కిస్తుంది;
అప్పుడు కొత్త స్థిరసంకేతాలు గల వాక్యాలకు గరిష్ఠత రాదు. కావాల్సిన
విస్తరించిన భాషలను పునరుద్ధరించాం.}
$n \ge 0$కు $(\Gamma_n, \Delta_n)$ అనే పెరుగుతున్న రెండు
\tetoken{వాక్యాల}{!!{sentence}s} సమితుల అనుక్రమాలను ఇలా నిర్వచిస్తాం.
$\Gamma_0 = \{ !A\}$, $\Delta_0 = \{\lnot !B \}$గా ఉంచండి.
$(\Gamma_n, \Delta_n)$ అప్పటికే నిర్వచించబడ్డాయని అనుకొని,
$\Gamma_{n+1}$, $\Delta_{n+1}$లను ఇలా నిర్వచించండి:
\begin{enumerate}
\item $\Gamma_n \cup \{!A_n \}$, $\Delta_n$లు $\Lang{L}'_0$లో
  వేరుపరచలేనివైతే $!A_n$ను $\Gamma_{n+1}$లో ఉంచండి. అంతేకాక $!A_n$
  అనేది $\lexists[x][!S]$ అనే అస్తిత్వాత్మక
  \tetoken{సూత్రం}{!!{formula}} అయితే, $\Gamma_n$, $\Delta_n$, $!A_n$,
  $!B_n$లలో రాని కొత్త \tetoken{స్థిరసంకేతం}{!!{constant}}~$c$ను
  ఎంచుకొని $\Subst{!S}{c}{x}$ను $\Gamma_{n+1}$లో ఉంచండి.
\item $\Gamma_{n+1}$, $\Delta_n \cup \{!B_n \}$లు $\Lang{L}'_0$లో
  వేరుపరచలేనివైతే $!B_n$ను $\Delta_{n+1}$లో ఉంచండి. అంతేకాక $!B_n$
  అనేది $\lexists[x][!S]$ అనే అస్తిత్వాత్మక
  \tetoken{సూత్రం}{!!{formula}} అయితే, $\Gamma_{n+1}$, $\Delta_n$,
  $!A_n$, $!B_n$లలో రాని కొత్త \tetoken{స్థిరసంకేతం}{!!{constant}}~$c$ను
  ఎంచుకొని $\Subst{!S}{c}{x}$ను $\Delta_{n+1}$లో ఉంచండి.
\end{enumerate}
ఇక్కడ పేర్కొన్న చేర్పు జరగని ప్రతి సందర్భంలో
$\Gamma_{n+1}=\Gamma_n$, $\Delta_{n+1}=\Delta_n$గా ఉంచాలి; చేర్పు
జరిగితే కొత్త సమితి ముందరి సమితిని కూడా పూర్తిగా కలిగి ఉంటుంది.
\sourcecorrection{OLTEMODINT-009}{ఆంగ్ల మూలం ఈ రెండు అనుక్రమాలు
పెరుగుతున్నవని చెప్పి, తరువాతి ఆగమనంలో చేర్పు జరగని సందర్భాలకు
$\Gamma_{n+1}=\Gamma_n$, $\Delta_{n+1}=\Delta_n$ను ఉపయోగిస్తుంది;
కానీ పునరావర్తక నిర్వచనంలో ఆ సందర్భాలను అసలు పేర్కొనదు. నిరూపణలో
అవసరమైన ఆ రెండు స్థిరీకరణ నిబంధనలను స్పష్టంగా చేర్చాం.}
చివరగా ఇలా నిర్వచించండి:
\begin{align*}
  \Gamma^* & = \bigcup_{n\ge 0} \Gamma_n, &
  \Delta^* & = \bigcup_{n\ge 0} \Delta_n.
\end{align*}
ఇప్పుడు $n$పై ఏకకాల ఆగమనం ద్వారా కింది రెండింటినీ నిరూపించవచ్చు:
\begin{enumerate}
\item\ollabel{part-a} $\Gamma_n$, $\Delta_n$లు $\Lang{L}'_0$లో
  వేరుపరచలేనివి;
\item\ollabel{part-b} $\Gamma_{n+1}$, $\Delta_n$లు $\Lang{L}'_0$లో
  వేరుపరచలేనివి.
\end{enumerate}
\olref{part-a} ఆధార సందర్భాన్ని \olref[sep]{lem:sep1} ఇస్తుంది.
\olref{part-b} అంశానికి మూడు సందర్భాలను వేరు చేయాలి:
\begin{enumerate}
\item $\Gamma_0 \cup \{!A_0 \}$, $\Delta_0$లు వేరుపరచదగినవైతే,
  $\Gamma_1 = \Gamma_0$; అప్పుడు \olref{part-b} కేవలం
  \olref{part-a}నే అవుతుంది;
\item $\Gamma_1 = \Gamma_0 \cup\{ !A_0\}$ అయితే, నిర్మాణం ప్రకారమే
  $\Gamma_1$, $\Delta_0$లు వేరుపరచలేనివి.
\item $!A_0$ అస్తిత్వాత్మకమైన మిగిలిన సందర్భంలో
  $\Gamma_1 = \Gamma_0 \cup \{ \lexists[x][!S],
  \Subst{!S}{c}{x} \}$. నిర్మాణం ప్రకారం
  $\Gamma_0 \cup \{ \lexists[x][!S]\}$, $\Delta_0$లు
  వేరుపరచలేనివి; అందువల్ల \olref[sep]{lem:sep2} ద్వారా
  $\Gamma_0 \cup \{ \lexists[x][!S], \Subst{!S}{c}{x} \}$,
  $\Delta_0$లు కూడా వేరుపరచలేనివే.
\end{enumerate}
దీనితో పై \olref{part-a}, \olref{part-b}ల ఆగమన ఆధారం పూర్తయింది.
ఇప్పుడు ఆగమన అడుగును చూద్దాం. \olref{part-a} కోసం,
$\Delta_{n+1} = \Delta_n \cup \{ !B_n \}$ అయితే నిర్మాణం ప్రకారం
$\Gamma_{n+1}$, $\Delta_{n+1}$లు వేరుపరచలేనివి ($!B_n$
అస్తిత్వాత్మకమైనప్పుడు కూడా \olref[sep]{lem:sep2} వల్ల ఇదే నిజం).
$\Gamma_{n+1}$, $\Delta_n \cup \{!B_n\}$లు వేరుపరచదగినవి కావడం వల్ల
$\Delta_{n+1} = \Delta_n$ అయితే, \olref{part-b}పై ఆగమన పరికల్పనను
వాడతాం. \olref{part-b} ఆగమన అడుగు కోసం,
$\Gamma_{n+2} = \Gamma_{n+1} \cup \{!A_{n+1} \}$ అయితే నిర్మాణం
ప్రకారం $\Gamma_{n+2}$, $\Delta_{n+1}$లు వేరుపరచలేనివి
($!A_{n+1}$ అస్తిత్వాత్మకమైనప్పుడు కూడా \olref[sep]{lem:sep2} వల్ల
ఇదే నిజం). $\Gamma_{n+2} = \Gamma_{n+1}$ అయితే, అప్పుడే నిరూపించిన
\olref{part-a} ఆగమన సందర్భాన్ని వాడతాం. దీనితో \olref{part-a},
\olref{part-b}లపై ఆగమనం పూర్తయింది.

దీనివల్ల $\Gamma^*$, $\Delta^*$లు వేరుపరచలేనివని వస్తుంది. లేకపోతే
సంహతత్వం వల్ల ఏదో $n \ge 0$ వద్ద $\Gamma_n$, $\Delta_n$లను
వేరుపరిచే వాక్యం ఉంటుంది; ఇది \olref{part-a}కు విరుద్ధం. ప్రత్యేకంగా,
$\Gamma^*$, $\Delta^*$లు అవైరుధ్యమైనవి: మొదటిది లేదా రెండవది
వైరుధ్యభరితమైతే వాటిని వరుసగా $\lexists[x][\eq/[x][x]]$ లేదా
$\lforall[x][\eq[x][x]]$ వేరుపరుస్తుంది.

ఇప్పుడు $\Lang{L}'_1$లో $\Gamma^*$ గరిష్ఠ అవైరుధ్య సమితి అని,
అలాగే $\Lang{L}'_2$లో $\Delta^*$ గరిష్ఠ అవైరుధ్య సమితి అని చూపుదాం.
మొదటిదాని కోసం, ఏదో $n \ge 0$కు $!A_n \notin \Gamma^*$,
$\lnot !A_n \notin \Gamma^*$ రెండూ అనుకుందాం. $!A_n \notin
\Gamma^*$ అయితే $\Gamma_n \cup \{!A_n \}$, $\Delta_n$ నుంచి
వేరుపరచదగినది. కాబట్టి కింది రెండూ తీరే $!C \in \Lang{L}'_0$
ఉంటుంది:
\begin{align*}
  \Gamma^* & \Entails !A_n \lif !C, &
  \Delta^* & \Entails \lnot !C.
\end{align*}
అలాగే $\lnot !A_n \notin \Gamma^*$ అయితే, కింది రెండూ తీరే
$!C' \in \Lang{L}'_0$ ఉంటుంది:
\begin{align*}
  \Gamma^* & \Entails \lnot !A_n \lif !C', &
  \Delta^* & \Entails \lnot !C'.
\end{align*}
ప్రతిజ్ఞా తర్కం ప్రకారం $\Gamma^* \Entails !C \lor !C'$,
$\Delta^* \Entails \lnot (!C \lor !C')$. అందువల్ల $!C \lor !C'$ అనేది
$\Gamma^*$, $\Delta^*$లను వేరుపరుస్తుంది. ఇలాంటి వాదనే $\Delta^*$
గరిష్ఠమని నిర్ధారిస్తుంది.

చివరగా $\Lang{L}'_0$లో $\Gamma^* \cap \Delta^*$ గరిష్ఠ అవైరుధ్య
సమితి అని చూపుదాం. అవైరుధ్య సమితుల ఛేదనం కనుక అది అవైరుధ్యమని
స్పష్టమే. గరిష్ఠత కోసం $!S \in \Lang{L}'_0$ను తీసుకోండి. ఇప్పుడు
$\Lang{L'_1} \supseteq \Lang{L'_0}$లో $\Gamma^*$ గరిష్ఠం; అలాగే
$\Lang{L'_2} \supseteq \Lang{L'_0}$లో $\Delta^*$ గరిష్ఠం. కనుక
$!S \in \Gamma^*$ లేదా $\lnot !S \in \Gamma^*$; అలాగే
$!S \in \Delta^*$ లేదా $\lnot !S \in \Delta^*$. ఒకవేళ
$!S \in \Gamma^*$, $\lnot !S \in \Delta^*$ అయితే $!S$ అనేది
$\Gamma^*$, $\Delta^*$లను వేరుపరుస్తుంది. ఒకవేళ
$\lnot !S \in \Gamma^*$, $!S \in \Delta^*$ అయితే $\lnot !S$ అనేది
$\Gamma^*$, $\Delta^*$లను వేరుపరుస్తుంది. అందువల్ల
$!S \in \Gamma^* \cap \Delta^*$ లేదా
$\lnot !S \in \Gamma^* \cap \Delta^*$; కాబట్టి
$\Gamma^* \cap \Delta^*$ గరిష్ఠం.

$\Gamma^*$ గరిష్ఠ అవైరుధ్య సమితి కనుక దానికి ఒక నమూనా
$\Struct{M}'_1$ ఉంటుంది. సంపూర్ణతా సిద్ధాంతపు నిరూపణలో ఉన్నట్లే
(\olref[fol][com][cth]{thm:completeness}), దాని
\tetoken{వ్యక్తి క్షేత్రం}{!!{domain}}~$\Domain{M'_1}$లో
\tetoken{స్థిరసంకేతాలకు}{!!{constant}s} అర్థనిర్దేశాలైన
$\Assign{c}{M'_1}$ మూలకాలే, అవన్నీ ఉంటాయి. అలాగే $\Delta^*$కు ఒక
నమూనా $\Struct{M}'_2$ ఉంటుంది; దాని
\tetoken{వ్యక్తి క్షేత్రం}{!!{domain}}~$\Domain{M'_2}$ను
\tetoken{స్థిరసంకేతాల}{!!{constant}s} అర్థనిర్దేశాలు
$\Assign{c}{M'_2}$ ఇస్తాయి.

$\Lang{L'_1} \setminus \Lang{L'_0}$లోని
\tetoken{స్థిరసంకేతాలు}{!!{constant}s},
\tetoken{ప్రమేయ సంకేతాలు}{!!{function}s},
\tetoken{విధేయ సంకేతాల}{!!{predicate}s} అర్థనిర్దేశాలను తొలగించి
$\Struct{M'_1}$ నుంచి $\Struct{M_1}$ను పొందండి; అలాగే
$\Struct{M_2}$ను పొందండి. $h(\Assign{c}{M'_1}) =
\Assign{c}{M'_2}$గా నిర్వచించిన $h \colon M_1 \to M_2$,
$\Lang{L}'_0$లో ఒక సమరూపత. ఎందుకంటే చూపినట్లుగా
$\Gamma^* \cap \Delta^*$, $\Lang{L}'_0$లో గరిష్ఠ అవైరుధ్య సమితి.
ప్రతి $\Lang{L}'_0$-
\tetoken{వాక్యం}{!!{sentence}} $\Gamma^*$, $\Delta^*$ రెండింటిలోనూ
ఉంటుంది లేదా రెండింటిలోనూ ఉండదు కాబట్టి ఇది అనుగమిస్తుంది: అందువల్ల
$\Assign{c}{M'_1} \in \Assign{P}{M'_1}$ అయినప్పుడు, అప్పుడు మాత్రమే
$\Atom{P}{c} \in \Gamma^*$; అయినప్పుడు, అప్పుడు మాత్రమే
$\Atom{P}{c} \in \Delta^*$; అయినప్పుడు, అప్పుడు మాత్రమే
$\Assign{c}{M'_2} \in \Assign{P}{M'_2}$. సమరూపతలు తీర్చాల్సిన ఇతర
షరతులనూ ఇలాగే నిర్ధారించవచ్చు.

ఇప్పుడు \tetoken{భాష}{!!{language}}~$\Lang{L_1} \cup \Lang{L_2}$కు
ఒక నమూనా $\Struct{M}$ను ఇలా నిర్వచిద్దాం:
\begin{enumerate}
\item \tetoken{వ్యక్తి క్షేత్రం}{!!{domain}}~$\Domain{M}$ కేవలం
  $\Domain{M_2}$, అంటే $\Assign{c}{M'_2}$ అనే అన్ని మూలకాల సమితి;
\item \tetoken{ఒక విధేయ సంకేతం}{!!a{predicate}}~$P$,
  $\Lang{L_2} \setminus \Lang{L_1}$లో ఉంటే
  $\Assign{P}{M} = \Assign{P}{M'_2}$;
\item విధేయ సంకేతం $P$, $\Lang{L}_1\setminus \Lang{L}_2$లో ఉంటే
  $\Assign{P}{M} = h(\Assign{P}{M'_1})$, అంటే
  $\tuple{\Assign{c_1}{M'_2}, \dots, \Assign{c_n}{M'_2}} \in
  \Assign{P}{M}$ అయినప్పుడు, అప్పుడు మాత్రమే
  $\tuple{\Assign{c_1}{M'_1}, \dots, \Assign{c_n}{M'_1}} \in
  \Assign{P}{M'_1}$.
  \sourcecorrection{OLTEMODINT-005}{$P$ అనేది
  $\Lang{L}_1\setminus\Lang{L}_2$కు చెందినదని అదే అంశం చెబుతుంది;
  కాబట్టి $\Struct{M'_2}$లో $P$కు అర్థనిర్దేశం ఉండదు. ఆంగ్ల మూలంలోని
  నిర్వచించని $h(\Assign{P}{M'_2})$ను, వెంటనే ఇచ్చిన ఘటకవారీ
  వివరణ కోరే $h(\Assign{P}{M'_1})$గా సరిచేశాం.}
\item \tetoken{ఒక విధేయ సంకేతం}{!!a{predicate}}~$P$,
  $\Lang{L}_0$లో ఉంటే $\Assign{P}{M} = \Assign{P}{M'_2} =
  h(\Assign{P}{M'_1})$.
\item $\Lang{L}_1 \cup \Lang{L}_2$లోని
  \tetoken{ప్రమేయ సంకేతాలను}{!!^{function}s},
  \tetoken{స్థిరసంకేతాలతో}{!!{constant}s} సహా, ఇలాగే నిర్వహించాలి.
\end{enumerate}
మూడు విస్తరించిన భాషలకూ పంచుకున్న కొత్త స్థిరసంకేతాల రెండవ హెన్కిన్
నిర్మాణపు అర్థనిర్దేశాలను కూడా ఇదే నిర్మాణంలో ఉంచి, చివరి రెండు
విస్తరించిన భాషల సమ్మేళనానికి గల విస్తరణగా పరిగణిస్తాం.
\sourcecorrection{OLTEMODINT-004}{ఆంగ్ల మూలం $\Struct{M}$ను
విస్తరించని $\Lang{L_1}\cup\Lang{L_2}$కు మాత్రమే నిర్వచించిన వెంటనే,
కొత్త స్థిరసంకేతాలు గల $\Gamma^*\cup\Delta^*$ను అది సత్యం చేస్తుందని
చెబుతుంది. ఆ వాక్యాల మదింపుకు అవసరమైన పంచుకున్న $c_n$ అర్థనిర్దేశాలను
నిలిపి ఉంచే విస్తరణను స్పష్టపరిచాం; మూల భాషకు సంకుచితం చేసినప్పుడు
చివరి $!A$, $\lnot!B$ ఫలితం మారదు.}

చివరగా, \tetoken{సూత్రాలపై}{!!{formula}s} ఆగమనం ద్వారా,
$\Struct{M}$ యొక్క సంతృప్తి $\Lang{L'_1}$లోని అన్ని
\tetoken{సూత్రాలకు}{!!{formula}s} పై సమరూపత ద్వారా రవాణా చేసిన చర
నిర్దేశాలపై $\Struct{M'_1}$ సంతృప్తితో ఒకటేనని; $\Lang{L'_2}$లోని
అన్ని \tetoken{సూత్రాలకు}{!!{formula}s} $\Struct{M'_2}$ సంతృప్తితో
ఒకటేనని చూపవచ్చు.
\sourcecorrection{OLTEMODINT-006}{రెండు నిర్మాణాల వ్యక్తి క్షేత్రాలు
వేరు కాబట్టి సూత్రాలపై అవి అక్షరార్థంగా “ఏకీభవిస్తాయి” అనడం సరిపోదు.
మొదటి వైపు చర నిర్దేశాన్ని సమరూపత $h$ ద్వారా రవాణా చేసినప్పుడు
సంతృప్తి సమానమనే అవసరమైన ఆగమన ప్రకటనను స్పష్టపరిచాం.}
ప్రత్యేకంగా $\Struct{M} \Entails \Gamma^* \cup \Delta^*$; అందువల్ల
$\Struct{M} \Entails !A$, $\Struct{M} \Entails \lnot!B$, ఇంకా
$\not\Entails !A \lif !B$. దీనితో క్రేగ్ అంతర్వేశ సిద్ధాంతపు
నిరూపణ పూర్తయింది.
\end{proof}

\end{document}
